\documentclass[11pt]{beamer}
\usetheme{Madrid}
\usecolortheme{seahorse}
\usepackage[utf8]{inputenc}
\usepackage[english]{babel}
\usepackage{amsmath,amssymb}
\usepackage{graphicx}
\usepackage{booktabs}
\usepackage{listings}
\usepackage{hyperref}
\usepackage{bbm}

\title[Systemic Risk Experiment]{Systemic Risk Analysis\\in Interbank Networks}
\subtitle{Detecting fragility at high connectivity}
\author{Interbank Model}
\institute{Interbank Simulation Framework}
\date{\today}

\begin{document}

\begin{frame}
\titlepage
\end{frame}

\begin{frame}{Motivation}
\centering
\Large
How does systemic risk behave\\as network connectivity increases?

\vspace{1em}
\normalsize
\begin{itemize}
    \item In sparsely connected networks, failures are isolated events
    \item In dense networks, ``mutual insurance'' can turn into a ``financial accelerator''
    \item Strategy synchronization amplifies shocks instead of absorbing them
\end{itemize}
\end{frame}

\begin{frame}{Research Question}
\centering
\Large
\textbf{Does high connectivity generate\\heavier tails in the bankruptcy distribution?}

\vspace{1.5em}
\normalsize
\noindent
If so:
\begin{itemize}
    \item Higher kurtosis $\Rightarrow$ extreme events more likely
    \item Lower Hill exponent $\Rightarrow$ heavier tail
    \item Lower distance to default $\Rightarrow$ higher leverage
\end{itemize}
\end{frame}

\begin{frame}{The Interbank Model}
\small
\noindent\textbf{Agents:} $N=50$ banks, each with balance sheet:
\[
C_i + L_i + R_i = D_i + E_i
\]

\medskip
\noindent\textbf{Time step} ($T=1000$ iterations):
\begin{enumerate}
    \item \textbf{Shock 1:} $\Delta D \sim \mathcal{U}(\mu, \mu + \omega)$,
          $d_i = \max(0, \Delta D_i - C_i)$, $s_i = \max(0, C_i - \Delta D_i)$
    \item \textbf{Network:} Erdős--Rényi graph with probability $p$
    \item \textbf{Interest rate:} $r_i = \dfrac{m}{\pi_i (1 - \psi_{\ell(i)})}$,
          where $\pi_i = E_i / \max E_{\text{borr}}$ and $\psi_j = E_j / \max E_{\text{lend}}$
    \item \textbf{Normalization (optional):} \texttt{robust2\_ir} — winsorization at percentiles 5--95
    \item \textbf{Loans:} $l_i = \min(\text{capacity}_i, s_{\ell(i)})$,
          rationing if $l_i < d_i$
    \item \textbf{Shock 2:} second deposit shock (only for borrowers)
    \item \textbf{Repayment:} loans and interests are settled;
          if $E_i < 0$ or $L_i < 0$ the bank fails
    \item \textbf{Replacement:} failed banks are re-initialized
\end{enumerate}
\end{frame}

\begin{frame}{robust2\_ir Normalization}
\framesubtitle{Winsorization without rescaling}

\noindent\textbf{Problem:} Raw interest rates may contain extreme values
that distort averages.

\vspace{0.5em}
\noindent\textbf{Solution:} Clip (winsorize) values outside a percentile range,
\emph{without rescaling}.

\vspace{0.5em}
\[
\tilde{r}_i =
\begin{cases}
Q_{\text{low}}  & \text{if } r_i < Q_{\text{low}} \\
r_i             & \text{if } Q_{\text{low}} \leq r_i \leq Q_{\text{high}} \\
Q_{\text{high}} & \text{if } r_i > Q_{\text{high}}
\end{cases}
\]

\vspace{0.5em}
\noindent With $Q_{\text{low}} = 5\%$, $Q_{\text{high}} = 95\%$ (configurable).

\medskip
\noindent\textbf{Properties:}
\begin{itemize}
    \item Monotonic $\Rightarrow$ preserves ir vs $\psi$ correlation
    \item No rescaling $\Rightarrow$ preserves natural trend with $p$
    \item Only removes outliers $\Rightarrow$ distribution shape almost intact
\end{itemize}
\end{frame}

\begin{frame}[fragile]{robust2\_ir Implementation}
\small
\begin{lstlisting}[language=Python]
elif self.config.robust2_ir:
    finite_mask = np.isfinite(self.interest_rate)
    if np.any(finite_mask):
        finite_values = self.interest_rate[finite_mask]
        q_low = float(self.config.robust_ir_quantile_low)    # 5
        q_high = float(self.config.robust_ir_quantile_high)  # 95
        low = np.percentile(finite_values, q_low)
        high = np.percentile(finite_values, q_high)
        if high > low:
            self.interest_rate[finite_mask] = np.clip(
                finite_values, low, high)
\end{lstlisting}

\vspace{0.3em}
\noindent Mutually exclusive with \texttt{normalize\_ir}, \texttt{sqrt\_ir}, \texttt{robust\_ir}.
\end{frame}

\begin{frame}{Experiment Design: \texttt{exp\_sysrisk.py}}
\small
\noindent\textbf{Fixed parameters:}
\begin{itemize}
    \item $N = 50$ banks, $T = 1000$ steps, $MC = 30$ Monte Carlo runs
    \item Bank replacement enabled
    \item Normalization: \texttt{robust2\_ir = True} (winsorization 5--95)
\end{itemize}

\medskip
\noindent\textbf{Independent variable:}
\[
p \in [0.08, 1.0] \quad \text{(20 evenly spaced values)}
\]

\medskip
\noindent\textbf{Additional metrics:}
\begin{itemize}
    \item \texttt{kurtosis\_equity}: Fisher kurtosis of the total equity series $E_t$
    \item \texttt{hill\_equity}: Hill tail exponent of the equity distribution
    \item \texttt{dist\_to\_default}: mean capital ratio $E/(D+E)$
\end{itemize}
\end{frame}

\begin{frame}{Metric 1: Kurtosis of Equity}
\centering
\Large
\[
\text{Kurtosis} = \frac{1}{n} \sum_{t=1}^T
\left( \frac{E_t - \bar{E}}{\sigma_E} \right)^4 - 3
\]

\vspace{0.5em}
\normalsize
\noindent
Where $E_t = \sum_i E_{i,t}$ is the total equity at each step $t$.

\vspace{1em}
\begin{columns}
\begin{column}{0.45\textwidth}
\centering
\textbf{Kurtosis $> 0$}\\
Heavier tail\\
Extreme equity\\
fluctuations
\end{column}
\begin{column}{0.45\textwidth}
\centering
\textbf{Kurtosis $< 0$}\\
Lighter tail\\
Equity stays near\\
the mean
\end{column}
\end{columns}

\vspace{0.5em}
\noindent\textbf{Hypothesis:} High $p$ $\Rightarrow$ higher kurtosis (more extreme equity swings).
\end{frame}

\begin{frame}{Metric 2: Hill Exponent of Equity}
\centering
\Large
\[
H_k = \frac{1}{k} \sum_{i=1}^{k} \log\left( \frac{E_{(i)}}{E_{(k+1)}} \right)
\]

\vspace{0.5em}
\normalsize
\noindent
$E_{(1)} \geq E_{(2)} \geq \cdots \geq E_{(k+1)}$ are the sorted
total equity values (upper tail = periods of peak capitalization).

\vspace{1em}
\begin{columns}
\begin{column}{0.45\textwidth}
\centering
\textbf{Low Hill}\\
Heavy tail\\
System spends more\\
time at peak equity\\
(extreme stability)
\end{column}
\begin{column}{0.45\textwidth}
\centering
\textbf{High Hill}\\
Light tail\\
Peak equity periods\\
are transient\\
(rapidly revert)
\end{column}
\end{columns}

\vspace{0.5em}
\noindent\textbf{Hypothesis:} High $p$ $\Rightarrow$ lower Hill exponent (longer periods of extreme capitalization).
\end{frame}

\begin{frame}{Metric 3: Distance to Default}
\centering
\Large
\[
\rho_t = \frac{E_t}{D_t + E_t}
\]

\vspace{0.5em}
\normalsize
\noindent
Aggregate capital ratio at each step $t$.

\vspace{1em}
\begin{itemize}
    \item $\rho$ high ($\approx 0.10$): well-capitalized banks
    \item $\rho$ low: leveraged banks, more fragile
    \item At equilibrium: $E_0 / (D_0 + E_0) = 1 / (9 + 1) = 0.10$
\end{itemize}

\vspace{0.5em}
\noindent\textbf{Hypothesis:} High $p$ $\Rightarrow$ lower $\rho$ (higher leverage
from synchronization of strategies).
\end{frame}

\begin{frame}{Running the Experiment}
\small
\noindent\textbf{From the command line:}
\begin{lstlisting}
python experiments/exp_sysrisk.py --do
\end{lstlisting}

\medskip
\noindent\textbf{Output:}
\begin{itemize}
    \item \texttt{output/exp\_sysrisk/results.gdt} — full results
    \item \texttt{output/exp\_sysrisk/*.png} — plots of each variable vs $p$
    \item \texttt{output/exp\_sysrisk/*.txt} — numerical data
    \item \texttt{output/exp\_sysrisk/results.txt} — cross-correlations
\end{itemize}

\medskip
\noindent\textbf{Clean and re-run:}
\begin{lstlisting}
python experiments/exp_sysrisk.py --do --clear
\end{lstlisting}
\end{frame}

\begin{frame}{Expected Interpretation}
\centering
\Large
If the dense network acts as a financial accelerator:

\vspace{1em}
\normalsize
\begin{columns}
\begin{column}{0.3\textwidth}
\centering
\includegraphics[width=\textwidth,height=4em]{example-image}\\
\textbf{Kurtosis of equity}\\
Increasing with $p$\\
(more extreme swings)
\end{column}
\begin{column}{0.3\textwidth}
\centering
\includegraphics[width=\textwidth,height=4em]{example-image}\\
\textbf{Hill of equity}\\
Decreasing with $p$\\
(longer peak periods)
\end{column}
\begin{column}{0.3\textwidth}
\centering
\includegraphics[width=\textwidth,height=4em]{example-image}\\
\textbf{Distance to default}\\
Decreasing with $p$\\
(higher leverage)
\end{column}
\end{columns}

\vspace{1em}
\noindent This would confirm that connectivity transforms
the network into a vector of systemic insolvency.
\end{frame}

\begin{frame}{Summary}
\centering
\Large
\textbf{Three equity-based metrics, one hypothesis}

\vspace{1.5em}
\normalsize
\begin{enumerate}
    \item \textbf{Equity kurtosis} $\uparrow$ $\Rightarrow$ more extreme equity fluctuations
    \item \textbf{Equity Hill exponent} $\downarrow$ $\Rightarrow$ longer periods of extreme capitalization
    \item \textbf{Capital ratio $\rho$} $\downarrow$ $\Rightarrow$ greater fragility
\end{enumerate}

\vspace{1.5em}
\noindent If all three point in the expected direction, we have evidence
that \textbf{high connectivity generates systemic fragility}.
\end{frame}

\begin{frame}{References}
\small
\begin{itemize}
    \item Allen, F. \& Gale, D. (2000). Financial contagion. \textit{Journal of Political Economy}, 108(1), 1--33.
    \item Battiston, S. et al. (2012). DebtRank: Too central to fail? \textit{Journal of Economic Dynamics and Control}, 36(11), 1615--1626.
    \item Hill, B. M. (1975). A simple general approach to inference about the tail of a distribution. \textit{The Annals of Statistics}, 3(5), 1163--1174.
    \item Source code: \url{https://github.com/anomalyco/interbank}
\end{itemize}
\end{frame}

\end{document}
